\documentclass[12pt]{article}
\usepackage[utf8]{inputenc}
\usepackage[T1]{fontenc}
\usepackage{hyperref}
\title{Foodo Language Resources for K-12 Curriculum Development}
\author{Compiled by Claude Sonnet 5.5 (claude-sonnet-5-5) for Andr\'as Kornai}
\date{}
\begin{document}
\maketitle

\section{Introduction}
Foodo (ISO 639-3 \texttt{fod}; Glottocode food1238) is a Guang language spoken around Sèmèrè in northern Benin. Search results give 21,000 native speakers (2018) in one source and about 37,000 people in another. Many speakers also use Tem with outsiders, which leads to confusion between the two languages.

\section{Primary Educational Resources}
None found. Omniglot says Foodo is written in the Latin alphabet but few speakers can read or write it.

\section{Linguistic and Reference Materials}
\begin{itemize}
\item Plunkett, ``An overview of Foodo, a linguistic island in Benin'', Journal of West African Languages: \url{https://journalofwestafricanlanguages.org/downloads/send/105-volume-36-number-1-2/540-an-overview-of-foodo-a-linguistic-island-in-benin}
\item Omniglot page on Foodo: \url{https://omniglot.com/writing/foodo.htm}
\item Wikipedia article: \url{https://en.wikipedia.org/wiki/Foodo_language}
\end{itemize}

\section{Religious and Community Resources}
Joshua Project lists Bible portions (1998) and no complete Bible or New Testament: \url{https://www.joshuaproject.net/index.php/people_groups/11223/BN}

\section{Digital and Audio Resources}
Omniglot notes that Foodo is heard on radio; no specific audio resource with a URL was found.

\section{Curriculum Development Context}
Foodo is not known to be taught in schools. Documentation is limited to a linguistic overview and scripture portions; a curriculum would need new primers and an agreed orthography.

\end{document}